% !TeX root = ../../Book3.tex
% 纲要标题：## 附录 G　实代数与半代数几何

\chapter{实代数与半代数几何}
\label{b3:app:G}

% 纲要细目（与计划纲要同步）：
% > 以第一卷附录 E 的有序域、实闭域为基础，结合本卷理想、坐标环和消去。范围是基本语言、代表性定理和小型证书，不展开完整实代数几何、模型论或优化理论。资料依据见修订说明 [R1]–[R3]。

\BookChapterNumber{b3:app:G}

\begin{chapterabstract}
实零点问题既需要理想，也需要多项式的符号。本附录从半代数投影出发，说明实根理想为何不同于普通根理想，继而比较非负性、多项式平方和与有理函数平方和。投影定理用迹形式判定一元符号，实零点定理与有理函数平方和则由有序域和证书推导。Sturm 序列给出精确实根计数和隔离，Gram 矩阵把二次型方法用于正性证书；紧集正性与矩表示给出构造概要，并注明所需分析工具。
\end{chapterabstract}

\begin{learninggoals}
\begin{itemize}
\item 区分实投影、消去理想所给的 Zariski 闭包，以及它们各自保留的信息。
\item 用平方和关系计算简单实根理想，并检查正性证书的适用范围。
\item 从有理系数多项式产生可核验的实根数与隔离区间。
\item 把平方和问题写成 Gram 矩阵条件，区别数值候选与精确证明。
\end{itemize}
\end{learninggoals}

所需基础为第一卷附录 E 的有序域和实闭域、第二卷第10.2节的实二次型，以及本卷第8章的 Gröbner 基与消去。紧集正性的概要另用一致估计，矩表示的概要还用 Stone--Weierstrass 与 Riesz 表示定理。前四节可按概念连续阅读，实零点定理调用的不可行性证书在第四节独立证明；若主要关心计算，可在看过半代数集合后先读最后一节的 Sturm 算法。Sturmfels《Solving Systems of Polynomial Equations》的 Section 1.3 与 Chapter 7 分别对应实根计算和平方和方法\cite[Section 1.3; Chapter 7]{Sturmfels2002PolynomialEquations}。

% 本附录深结果的证明状态分别在正文标明。

% !TeX root = ../../Book3.tex
% 纲要标题：### G.1 半代数集合与量词消去

\section{半代数集合与量词消去}
\label{b3:sec:G01}

实零点会随系数的符号变化，仅有多项式等式不足以描述参数条件。把等式与不等式一起保存在半代数集合中，消去变量后仍可得到有限的符号描述。

% 纲要细目（与计划纲要同步）：
% * 实闭域上由多项式等式与不等式及有限布尔运算定义的集合；半代数映射
% * Tarski–Seidenberg 投影定理与实闭域量词消去，说明它们为何超出普通 Zariski 消去
% * 手算一元或二元投影例子；完整柱状代数分解算法不放入正文
%
% #### G.1.1 半代数集合与半代数映射
% #### G.1.2 Tarski–Seidenberg 投影定理与量词消去
% #### G.1.3 低维投影算例与柱状代数分解


% 正文按纲要完成；所列定理给出证明或证明概要。

\subsection{半代数集合与半代数映射}
\label{b3:subsec:G01:01}

实系数方程 \(x^2+1=0\) 没有实解，而 \(x^2-a=0\) 是否有实解取决于参数 \(a\) 的符号。这说明研究实零点时，等式与不等式必须一起考虑。本附录用 \(R\) 表示实闭域；涉及紧性、光滑性和测度时，另明确限定为 \(\mathbb R\)。第一卷附录 E 已介绍实闭域及多项式的介值性质。

\begin{definition}\label{b3:def:semialgebraic-set}
设 \(R\) 为实闭域，\(m,n\) 为非负整数。若 \(R^n\) 的一个子集可由有限多个条件 \(f(x)=0\)、\(g(x)>0\)（其中 \(f,g\in R[x_1,\ldots,x_n]\)）经有限次并、交、补得到，则称它为\term{半代数集合}[semialgebraic set]。若 \(A\subseteq R^m\)、\(B\subseteq R^n\) 都是半代数集合，且映射 \(\varphi:A\to B\) 的图为 \(R^{m+n}\) 的半代数子集，则称 \(\varphi\) 为\term{半代数映射}[semialgebraic map]。
\end{definition}

条件 \(g\ge0\) 等价于 \(g>0\) 或 \(g=0\)，故定义允许弱不等式。把有限个多项式的符号取值逐一列出，还可把任意半代数集合写成有限个由等式与严格不等式定义的集合之并。多项式映射是半代数映射；有理函数只在其分母非零的定义域上给出半代数映射。例如 \(x\mapsto1/x\) 的图像由 \(xy=1\) 定义。一般半代数映射不必连续：符号函数就是例子。上述用语可参阅 Mangolte《Real Algebraic Varieties》\cite[Appendix B, Definitions B.2.1--B.2.2]{Mangolte2020RealAlgebraicVarieties}。

\subsection{Tarski–Seidenberg 投影定理与量词消去}
\label{b3:subsec:G01:02}

\begin{theorem}[Tarski--Seidenberg 投影定理]\label{b3:thm:real-projection}
设 \(R\) 为实闭域，\(S\subseteq R^{n+m}\) 为半代数集合。坐标投影 \(\pi:R^{n+m}\to R^n\) 的像 \(\pi(S)\) 仍为半代数集合。
\end{theorem}

\begin{proof}
逐个投影，归约为消去一个变量 \(T\)。将定义 \(S\) 的布尔公式中出现的多项式记为
\[
p_1(x,T),\ldots,p_l(x,T).
\]
对一个固定参数 \(x\)，我们要判定哪几个符号元组能由某个 \(T\in R\) 实现；若这一判定只需有限个参数多项式的符号，投影便为半代数集。

先按各 \(p_i\) 的 \(T\) 次数分支，次数由系数的零与非零条件确定；恒零多项式单独处理。用 Euclid 算法从其余 \(p_i\) 的乘积取得首一平方自由部分 \(P\)。每次除法按当前首项系数是否为零分支，次数下降保证只有有限步、有限个分支；各分支上 \(P\) 的系数是参数的有理函数，所用分母已限定非零。若 \(P=1\)，所有 \(p_i\) 在整条直线上符号恒定，直接由首项系数判定。以下设 \(\deg P=d>0\)。

令 \(B=R[T]/(P)\)。对于任意 \(q\in R[T]\)，在基 \(1,T,\ldots,T^{d-1}\) 上计算迹二次型
\[
H_q(u,v)=\operatorname{Tr}_{B/R}(quv).
\]
它的矩阵只用 \(P,q\) 的系数作有限次域运算。实闭域上，平方自由代数 \(B\) 分解成若干 \(R\) 与 \(R(\mathrm i)\) 的直积。在实根 \(a\) 对应的分量中，二次型为 \(q(a)u^2\)，符号差为 \(\operatorname{sign}q(a)\)；在非实共轭对对应的分量中，非零迹形式有一正一负，\(q\) 为零时形式也为零。因此
\[
\operatorname{signature}(H_q)
=\sum_{P(a)=0,\ a\in R}\operatorname{sign}q(a).
\]
这里符号差指正惯性指数减负惯性指数。对称矩阵可按非零对角元作合同消元；若全部对角元为零但某个非对角元非零，就取对应的二阶块作为枢轴。按枢轴的零与符号分支，有限步可算出惯性指数。所以迹形式的符号差由有限多个参数有理函数的符号决定，而分母非零时这些仍能写成多项式符号条件。迹形式的这一根计数原理也见\cite[Section 2.4, Theorem 2.8]{Sturmfels2002PolynomialEquations}。

把所有 \(p_i\) 及其 \(T\) 导数列为 \(q_1,\ldots,q_b\)。对每个 \(e\in\{0,1,2\}^b\)，计算 \(H_{\prod_jq_j^{e_j}}\) 的符号差。若 \(N_\sigma\) 表示这些 \(q_j\) 在 \(P\) 的实根上取符号元组 \(\sigma\in\{-1,0,1\}^b\) 的根数，则这些差给出有限线性方程
\[
Q_e=\sum_\sigma N_\sigma\prod_j\sigma_j^{e_j},
\qquad 0^0=1.
\]
一变量的系数矩阵是 \(1,s,s^2\) 在 \(s=-1,0,1\) 上的取值矩阵，行列式非零；多变量矩阵是它们的张量积，故也可逆。于是知道各 \(Q_e\) 就能确定每个 \(N_\sigma\)，从而判定哪些根上的符号元组实际出现。

这些数据也确定各根之间开区间上的符号元组。在一个根 \(a\) 的右边，\(p_i\) 的符号由其在 \(a\) 处首个非零导数的符号给出；这个符号在到下一根之前不变。最左边区间的符号由 \((-1)^{\deg p_i}\) 乘首项系数的符号给出。有限实根之间均有元素，实闭域的多项式介值性质保证开区间内没有未计入的变号。因此根上的元组、各根右侧的元组与最左侧元组已经穷尽所有 \(T\) 的可能性，无须先求出根的位置。

检查这些有限元组是否满足原布尔公式，就判定了 \(\exists T\)。前述所有分支均由有限个参数多项式的零、正、负条件给出；合并满足条件的分支，得到投影的半代数描述。反复消去 \(m\) 个变量即得结论。
\end{proof}

投影性质的概述见\cite[Chapter 2, p.~23]{BochnakCosteRoy1998RealAG}。条件
\[
\exists y_1,\ldots,y_m\quad \Phi(x,y)
\]
描述的正是投影，其中 \(\Phi\) 是多项式符号条件的有限布尔组合。定理把它改写成只含 \(x\) 的有限布尔组合；全称量词则用 \(\forall y\,\Phi\iff\neg\exists y\,\neg\Phi\) 处理。逐层消去量词，便得到实闭域理论的量词消去。

\ProseParagraphBreak
这与\cref{b3:thm:elimination-projection-closure}的任务不同。该定理在代数闭域上用消去理想描述投影的 Zariski 闭包，不能单凭等式记住实数的符号。例如 \(I=(y^2-x)\subseteq R[x,y]\) 满足 \(I\cap R[x]=0\)，但实投影是 \(\{x\ge0\}\)，并非整个 \(R\)。实代数消去要保留的是实际投影，因此输出自然包含不等式。

\subsection{低维投影算例与柱状代数分解}
\label{b3:subsec:G01:03}

\begin{example}\label{b3:exm:real-projection-disk}
设 \(R\) 为实闭域。集合
\[
S=\{(x,y):y^2+x^2=1,\ y\ge x\}
\]
在 \(x\) 轴上的投影为 \(\{-1\le x\le0\}\cup\{0\le x,\ 2x^2\le1\}\)。事实上，圆上有点首先要求 \(-1\le x\le1\)。若 \(x\le0\)，取非负根 \(y=\sqrt{1-x^2}\) 即可；若 \(x>0\)，可取的最大 \(y\) 为该非负根，故存在满足条件的 \(y\) 当且仅当 \(1-x^2\ge x^2\)。这一计算明确使用了实闭域中非负元有平方根，而非仅使用复数域上的消去。
\end{example}

\term{柱状代数分解}[cylindrical algebraic decomposition, CAD] 把实空间分成有限个半代数胞腔，使指定多项式在每个胞腔上符号不变，并使各胞腔向低维坐标空间的投影相同或不相交。在二维，先把 \(x\) 轴按若干关键根分成点与区间；在各区间上排列 \(y\) 方向的根图像，再取图像之间的带状区域。这正是“柱状”的含义，并非说各胞腔都是直积。

\ProseParagraphBreak
上例的关键参数包括 \(\pm1\) 以及 \(2x^2-1\) 的根。一般计算还须处理首项系数消失、重根与两个多项式共有根等情形；判别式与结式提供候选投影多项式，但退化时不能只凭一份通用公式完成分解。\secref{b3:sec:G05}给出一维根计数和隔离的完整步骤；高维 CAD 的正确性与复杂性证明留待实代数几何专门课程。

% !TeX root = ../../Book3.tex
% 纲要标题：### G.2 实根理想与实零点定理
% 纲要位置：计划纲要.md，第 2419 行。

\section{实根理想与实零点定理}
\label{b3:sec:G02}

% TODO: 按以下纲要要点撰写本节。

% 纲要内容（仅作写作计划，尚非教材正文）：
%
% * 实理想、实根理想和实零点集合；区分普通根理想与实根理想
% * 以 \(x^2+y^2\) 在 \(\mathbb R^2\) 上的零点说明普通根式不足以描述实消失理想
% * 实零点定理陈述与用途；一般环的实谱以素理想连同剩余分式域上的序作为入门对象，不展开完整谱理论
%

\subsection{实理想、实根理想与实零点集}
\label{b3:subsec:G02:01}

待撰写。

\subsection{平方和方程的实消失理想}
\label{b3:subsec:G02:02}

待撰写。

\subsection{实零点定理与实谱的基本对象}
\label{b3:subsec:G02:03}

待撰写。

% !TeX root = ../../Book3.tex
% 纲要标题：### G.3 非负性、平方和与 Hilbert 第17问题
% 纲要位置：计划纲要.md，第 2425 行。

\section{非负性、平方和与 Hilbert 第17问题}
\label{b3:sec:G03}

% TODO: 按以下纲要要点撰写本节。

% 纲要内容（仅作写作计划，尚非教材正文）：
%
% * 非负多项式、多项式平方和与有理函数平方和是三个需要区分的条件
% * 以 Motzkin 多项式说明非负性不保证多项式平方和；Artin 定理的有理函数平方和表述
% * 指明系数域、分母和定义域；仅选代表性证明或证明框架，不展开一般有效次数界
%

待撰写。

% !TeX root = ../../Book3.tex
% 纲要标题：### G.4 正性证书与 Positivstellensatz

\section{正性证书与 Positivstellensatz}
\label{b3:sec:G04}

约束集合上的正性可以通过约束多项式参与的平方和恒等式来验证。是否允许约束间的乘积、是否具有紧性或 Archimedes 条件，会决定证书的形式与存在范围。

% 纲要细目（与计划纲要同步）：
% * 预序、二次模及多项式不等式系统；以恒等式作为不可行性或正性的证书
% * Krivine–Stengle 型 Positivstellensatz 的一个明确版本；等式约束、严格不等式和非严格不等式分别处理
% * Schmüdgen／Putinar 型结果只作比较：紧性、Archimedean 条件及严格正性不能无条件互换
%
% #### G.4.1 预序、二次模与多项式不等式证书
% #### G.4.2 Krivine–Stengle 型 Positivstellensatz
% #### G.4.3 Schmüdgen 与 Putinar 型结果的假设比较


% 正文按纲要完成；所列定理给出证明或证明概要。

\subsection{预序、二次模与多项式不等式证书}
\label{b3:subsec:G04:01}

设约束集合为 \(K=\{x\in\mathbb R^n:g_1(x)\ge0,\ldots,g_s(x)\ge0\}\)。若能写出 \(f=\sigma_0+\sum_i\sigma_i g_i\)，各 \(\sigma_i\) 都为平方和，就已经证明 \(f\) 在 \(K\) 上非负。再允许约束之间的乘积，会得到更广的一类证书。

\begin{definition}\label{b3:def:preordering-quadratic-module}
设 \(R\) 为实闭域，\(A=R[x_1,\ldots,x_n]\)，记 \(\Sigma A^2\) 为有限平方和的集合。给定 \(g_1,\ldots,g_s\in A\)，集合
\[
M(g)=\Sigma A^2+\sum_{i=1}^s\Sigma A^2g_i,
\qquad
T(g)=\sum_{\varepsilon\in\{0,1\}^s}\Sigma A^2g_1^{\varepsilon_1}\cdots g_s^{\varepsilon_s}
\]
分别称为它们生成的\term{二次模}[quadratic module]与\term{预序}[preordering]。
\end{definition}

这里的二次模是含 \(1\)、对加法与乘平方封闭的子集，不是通常意义的 \(A\) 模。预序还对乘法封闭：把乘积中 \(g_i\) 的偶次幂吸收到平方和系数即可。二者中每个元素都在 \(K\) 上非负。比如在 \([-1,1]\) 上，\(1-x=(1-x)^2/2+(1-x^2)/2\) 是由约束 \(1-x^2\ge0\) 给出的精确证书。

\subsection{Krivine–Stengle 型 Positivstellensatz}
\label{b3:subsec:G04:02}

\begin{theorem}[Krivine--Stengle 型不可行性证书]\label{b3:thm:real-infeasibility-certificate}
设 \(R\) 为实闭域，\(f_1,\ldots,f_r,g_1,\ldots,g_s\in R[x_1,\ldots,x_n]\)，并令 \(I=(f_1,\ldots,f_r)\)、\(T(g)\) 为 \(g_1,\ldots,g_s\) 生成的预序。则
\[
\{x\in R^n:f_i(x)=0\ (1\le i\le r),\ g_j(x)\ge0\ (1\le j\le s)\}=\varnothing
\quad\Longleftrightarrow\quad -1\in I+T(g).
\]
\end{theorem}

\begin{proof}
若 \(-1\in I+T(g)\)，在任何假定的解上求值，理想部分为零、预序部分非负，不可能得到 \(-1\)。证明反向的逆否命题。设 \(-1\notin I+T(g)\)，在 \(B=R[x]/I\) 中取由 \(g_j\) 的像生成的预序，利用 Zorn 引理把它扩为不含 \(-1\) 的极大预序 \(P\)。预序包含全部平方、对加法与乘法封闭；链的并仍有这些性质且不含 \(-1\)。

说明极大性使它成为一个带素支集的正锥。若 \(a,-a\) 都不在 \(P\)，加入它们中的任一个所得预序分别为 \(P+aP\)、\(P-aP\)，因为 \(a^2\in P\)。两者都不能再扩大 \(P\)，于是有
\(-1=u+av\)、\(-1=w-az\)，其中 \(u,v,w,z\in P\)。相乘整理得
\(1+u+w+uw+a^2vz=0\)，即 \(-1\in P\)，矛盾。因此 \(P\cup(-P)=B\)。集合 \(\mathfrak q=P\cap(-P)\) 是理想：乘任意元时，可根据该元或其负元属于 \(P\) 来验证封闭性。

它还是素理想。若 \(ab\in\mathfrak q\) 而 \(a,b\notin\mathfrak q\)，更换符号后可令 \(a,b\in P\)、\(-a,-b\notin P\)。极大性给出 \(-1=u-av\)、\(-1=w-bz\)，故
\(abvz=(1+u)(1+w)\)。左边属于 \(\mathfrak q\)，把它的负元与 \(u+w+uw\in P\) 相加，再次得到 \(-1\in P\)，矛盾。

于是 \(B/\mathfrak q\) 是有序整环，其分式域 \(L\) 带有延拓的序，且各 \(g_j\) 的像非负。\(R\) 在 \(L\) 中仍保持原序，因为实闭域的正元都是非零平方。取 \(L\) 的一个实闭包 \(L^{\mathrm{rc}}\)，\(x_i\) 的像在这个实闭扩域内满足原系统。\cref{b3:thm:real-projection}给出的量词消去说明带 \(R\) 系数的存在公式在 \(R\subseteq L^{\mathrm{rc}}\) 中真值相同：消去后只剩 \(R\) 中元素的符号，而嵌入保持这些符号。因此原系统在 \(R\) 中也有解。
\end{proof}

实闭包的构造见第一卷附录 E。所证版本在 \(R=\mathbb R\) 时是\cite[Section 7.4, Theorem 7.5]{Sturmfels2002PolynomialEquations}没有严格不等式的情形；证明没有调用实零点定理，所以稍前使用此结果推导实根理想公式不会产生循环。

\ProseParagraphBreak
严格不等式和不等条件也可以转换后使用该定理。例如 \(h(x)\ne0\) 等价于存在 \(z\) 使 \(zh(x)-1=0\)；\(h(x)>0\) 等价于存在 \(z\) 使 \(h(x)z^2-1=0\)。后一等价使用正实数有平方根。转换增加变量，却保留了精确可行性。

\begin{example}\label{b3:exm:inconsistent-inequalities}
条件 \(x\ge0\)、\(-x-1\ge0\) 不可能同时成立，恒等式 \(-1=x+(-x-1)\) 就是预序证书。若另有等式 \(y-x^2=0\) 并要求 \(y<0\)，直接利用 \(y=x^2\) 已知矛盾；上述变量替换还给出等式证书：加入 \((-y)z^2-1=0\) 后，
\[
-1=(-yz^2-1)+z^2(y-x^2)+(xz)^2.
\]
最后一项为平方，前两项属于等式理想。
\end{example}

\subsection{Schmüdgen 与 Putinar 型结果的假设比较}
\label{b3:subsec:G04:03}

\begin{theorem}[紧集上的严格正性]\label{b3:thm:compact-positivity-certificates}
设 \(g_1,\ldots,g_s,f\in\mathbb R[x_1,\ldots,x_n]\)，\(K=\{x\in\mathbb R^n:g_i(x)\ge0\text{ 对每个 }i\}\)，并记 \(T(g)\)、\(M(g)\) 分别为 \(g_1,\ldots,g_s\) 生成的预序与二次模。
\begin{enumerate}
\item 若 \(K\) 紧且 \(f>0\) 于 \(K\)，则 \(f\in T(g)\)；这是 Schmüdgen 型结论。
\item 若存在 \(N>0\) 使 \(N-\sum_i x_i^2\in M(g)\)，且 \(f>0\) 于 \(K\)，则 \(f\in M(g)\)；这是 Putinar 型结论。
\end{enumerate}
\end{theorem}

\begin{proofsketch}
先证第二项，采用\cite[Section 2, Theorem 6, Lemma 8 and proof of Theorem 3]{Schweighofer2005Optimization}的构造。以下置 \(A=\mathbb R[x_1,\ldots,x_n]\)。记 \(b=N-\sum_i x_i^2\in M(g)\)，\(a=N+1/4\)，并令
\[
\Delta=\{y\in\mathbb R_{\ge0}^{2n}:\textstyle\sum_jy_j=2na\},
\qquad \ell(y)_i=(y_i-y_{n+i})/2,
\qquad C=\ell(\Delta).
\]
这是紧的单纯形及其线性像。把各 \(g_i\) 乘以适当正常数，使它们在 \(C\) 上都不超过一；这不改变 \(K\) 或 \(M(g)\)。可以选择 \(c\ge0\) 及充分大的整数 \(k\)，使
\[
q=f-c\sum_i(1-g_i)^{2k}g_i>0\quad\text{于 }C.
\]
理由是：在紧集 \(\{x\in C:f(x)\le\varepsilon\}\) 上，取 \(\varepsilon>0\) 使它与 \(K\) 不交，则至少有一个 \(g_i\le-\delta\)，其中 \(\delta>0\) 可统一选择；这个负项的扣除可使 \(f\) 增大。对于 \(0\le t\le1\)，有 \(t(1-t)^{2k}\le1/(2k+1)\)，所以所有非负项造成的损失随 \(k\) 增大而统一趋于零。先选 \(c\)，再选 \(k\)，便使上式在这个坏集及其补集都为正；坏集为空时取 \(c=0\)。

将 \(q\) 经 \(\ell\) 提升到 \(\Delta\)，再用 \(\sum y_j/(2na)\) 把各项补齐次数，得到在非负象限除原点以外严格正的齐次式 \(Q(y)\)。P\'olya 的系数引理保证，对充分大的 \(e\)，
\(Q(y)(\sum y_j/(2na))^e\) 的系数全非负。引理的证明是比较其按多项式系数归一化后的系数与 \(Q(\alpha/|\alpha|)\)：前者中的下降阶乘与普通幂之差为一致的 \(O(1/e)\)，而 \(Q\) 在标准单纯形上的最小值严格大于零。本概要省略这个有限系数估计的展开，详见上述 Theorem 6。

代入 \(y_i=a+x_i\)、\(y_{n+i}=a-x_i\) 后，\(\sum y_j/(2na)=1\)，所得式正是 \(q\)。而
\[
a\pm x_i=b+\sum_{j\ne i}x_j^2+(x_i\pm1/2)^2
\in T(b)=\Sigma A^2+b\Sigma A^2\subseteq M(g).
\]
\(T(b)\) 对乘法封闭，故上述非负系数表示给出 \(q\in T(b)\)。此前减去的每项属于 \(M(g)\)，所以 \(f\in M(g)\)。这里没有把一般二次模当作对乘法封闭的预序。

第一项还须从紧性得到 \(T(g)\) 内的球约束。以下把这个步骤写出。记 \(\Sigma=\sum_i x_i^2\)，选 \(s>0\) 使 \(h=s-\Sigma>0\) 于 \(K\)。系统 \(g_i\ge0,-h\ge0\) 不可行；\cref{b3:thm:real-infeasibility-certificate}给出 \(h t_1=1+t_0\)，其中 \(t_0,t_1\in T(g)\)。由恒等式
\[
h\bigl(1+(h-2)^2t_1\bigr)
=1+(h-3/2)^2+3/4+(h-2)^2t_0
\]
得到 \(h(1+v)=1+u\)，其中 \(u,v\in T(g)\)。令 \(T(h)=\Sigma A^2+h\Sigma A^2\)。对任意多项式 \(p\)，均存在常数 \(t\ge0\) 使 \(t\pm p\in T(h)\)：具有此性质的多项式对加法、乘法封闭，乘法用
\[
tu\pm pq=\tfrac12\bigl((t+p)(u\pm q)+(t-p)(u\mp q)\bigr)
\]
验证；而 \(s+1/4\pm x_i=h+\sum_{j\ne i}x_j^2+(x_i\pm1/2)^2\) 说明它包含各变量。

因此可取 \(t\ge0\) 使 \(t-sv\in T(h)\)。由于 \(1+v\in T(g)\) 且 \(h(1+v)\in T(g)\)，有
\((1+v)(t-sv)\in T(g)\)。另外 \(h+sv=h(1+v)+v\Sigma\in T(g)\)。将这两式与平方 \(s(t/(2s)-v)^2\) 相加，得到
\[
\bigl(s+t+t^2/(4s)\bigr)-\Sigma\in T(g).
\]
于是把已证第二项用于二次模 \(T(g)\)，即得 \(f\in T(g)\)。此球约束构造也说明为什么这一步依赖预序的乘法封闭性。所用弱紧性结论及两项的关系可对照\cite[Section 1, Theorems 1 and 3]{Schweighofer2005Optimization}。
\end{proofsketch}

第二个条件称为二次模的 Archimedean 条件的一种等价形式，它显式给出球约束，因而蕴含 \(K\) 紧；仅知 \(K\) 紧不能直接删去这个条件。两个结论的证书格式也不同：预序允许 \(g_ig_j\)，二次模只允许单个生成约束。

\ProseParagraphBreak
严格正性不能任意换成非负性。例如令 \(g=(1-x^2)^3\)，则 \(K=[-1,1]\)，\(M(g)\) 满足上述条件：由定理第一条及只有一个生成元时 \(T(g)=M(g)\)，可知 \(2-x^2\in M(g)\)。然而 \(1-x^2\notin M(g)\)。若有 \(1-x^2=\sigma_0+\sigma_1(1-x^2)^3\)，在 \(x=1\) 得 \(\sigma_0(1)=0\)，于是组成 \(\sigma_0\) 的每个平方的底数在 \(1\) 为零，故 \(\sigma_0'(1)=0\)。右边导数在 \(1\) 为零，左边却为 \(-2\)。这个简单的消失阶数检验解释了为何定理要保留 \(f>0\)。

% !TeX root = ../../Book3.tex
% 纲要标题：### G.5 实根计算、平方和优化与相关专题

\section{实根计算、平方和优化与相关专题}
\label{b3:sec:G05}

近似数值可以帮助寻找实根或平方和表示，最终核验仍要依靠精确的符号与矩阵数据。根计数给出隔离区间的证明，Gram 矩阵则把平方和问题转为带半正定条件的有限线性关系。

% 纲要细目（与计划纲要同步）：
% * Sturm 序列、实根隔离和柱状代数分解的任务；有效算法注明精确系数与比较运算的输入条件
% * 平方和的 Gram 矩阵、半正定规划与矩问题仅作简要介绍，区别浮点近似和精确代数证书
% * 半代数拓扑、Nash 几何、o-minimal 几何及完整矩问题转相应专题；不要求在学习交换代数时掌握这些领域的完整理论
%
% #### G.5.1 Sturm 序列、实根隔离与柱状代数分解
% #### G.5.2 Gram 矩阵、半正定规划与精确代数证书
% #### G.5.3 半代数拓扑、Nash 几何、o-minimal 几何与矩问题
%
% ---


% 正文按纲要完成；所列定理给出证明或证明概要。

\subsection{Sturm 序列、实根隔离与柱状代数分解}
\label{b3:subsec:G05:01}

\begin{theorem}[Sturm 根计数]\label{b3:thm:appendix-sturm-count}
设 \(p\in\mathbb Q[x]\) 非常数且无重因子，置 \(p_0=p,p_1=p'\)，递归作带负号的余式
\[
p_{j-1}=q_jp_j-p_{j+1},
\]
直到末项为非零常数。对实数 \(t\)，删去 \((p_0(t),\ldots,p_m(t))\) 中的零项后，记符号变化次数为 \(V(t)\)。若实数 \(a<b\) 满足 \(p(a)p(b)\ne0\)，则 \((a,b)\) 内的实根个数为 \(V(a)-V(b)\)。
\end{theorem}
\begin{proof}
Euclid 算法与 \(\gcd(p,p')=1\) 保证末项为常数，且相邻项没有公共零点。在一个中间项 \(p_j\) 的零点处，余式关系给出 \(p_{j-1}=-p_{j+1}\ne0\)。因此无论 \(p_j\) 从哪种符号变为哪种符号，它与两侧项合计恰有一次符号变化，变化总数不跳动。若多个不相邻中间项同时为零，各自控制互不重复的相邻符号对，结论相同。只有 \(p_0\) 的根可能改变 \(V\)。在单根 \(c\) 附近，\(p_0(x)\) 的符号等于 \((x-c)p_1(c)\) 的符号，而 \(p_1(x)\) 符号不变；穿过 \(c\) 时第一对由异号变同号，使 \(V\) 减一。各项只有有限多个根，分段相加即得公式。
\end{proof}

这也回顾了第一卷附录 E.3 的 Sturm 定理；计算表述见\cite[Section 1.3, Theorem 1.4]{Sturmfels2002PolynomialEquations}。有重根时先以 \(p/\gcd(p,p')\) 代替 \(p\)，从而计数不同实根；若还要重数，另保留平方自由分解。

\begin{example}\label{b3:exm:sturm-quartic-isolation}
对 \(p=x^4-5x^2+4\)，可把各非零余式乘以正有理常数简化，得到符号等价的 Sturm 序列
\[
x^4-5x^2+4,\quad4x^3-10x,\quad5x^2-8,\quad x,\quad1.
\]
在 \(-3,-3/2,0,3/2,3\) 处的符号分别为
\[
\begin{array}{c|ccccc|c}
t&p_0&p_1&p_2&p_3&p_4&V(t)\\\hline
-3&+&-&+&-&+&4\\
-3/2&-&+&+&-&+&3\\
0&+&0&-&0&+&2\\
3/2&-&-&+&+&+&1\\
3&+&+&+&+&+&0
\end{array}
\]
故相邻五个端点给出的四个区间各含一个实根，其他地方没有实根。虽然本例可直接分解为 \((x^2-1)(x^2-4)\)，计数过程只用有理数运算和符号比较。
\end{example}

\begin{proposition}[有理系数实根隔离]\label{b3:prop:rational-real-root-isolation}
给定非零 \(p\in\mathbb Q[x]\)，可用有限次有理数运算及整数比较，输出所有不同实根的两两不交有理开区间，每个区间恰含一根；有理根也可单列为精确数值。
\end{proposition}
\begin{proof}
若 \(p\) 为非零常数，返回空表。否则先去重因子并化成首一。若 \(p=x^d+\sum_{i<d}a_ix^i\)，取整数 \(B>1+\max_i|a_i|\)。当 \(|x|\ge B\) 时，低次项绝对值之和小于 \(|x|^d\)，所以全部实根在 \((-B,B)\)。以 Sturm 公式计算区间根数：零根区间丢弃，一根区间保留，多根区间继续分割。

分割一个区间 \((a,b)\) 时，在其中间三分之一内依次试取有理数，选到一个满足 \(p(c)\ne0\) 的数 \(c\)，再分为 \((a,c)\)、\((c,b)\)。非零多项式只有有限个根，故选点总会结束；所有区间端点也始终不是根。每次分割后子区间长度至多为原来的 \(2/3\)。若至少有两个不同实根，它们之间存在最小正间距；区间长度小于该间距后不可能仍含多根，所以全部分割在有限步内结束。只有零个或一个实根时，第一次计数便已结束。最终保留的区间两两不交，各含一根，且涵盖全部实根。算法遇到有理根时可以将其单列，但并不需要为此删除因子或更换正在计数的多项式。
\end{proof}

精确实代数系数也可用相应的域运算及序比较处理，但必须提供这些操作的有效表示。任意实数只给若干小数位时，无法据此保证重根判断和终止性。高维 CAD 反复需要一维隔离，却还增加参数退化分析；不能把上述一维终止性当作完整 CAD 正确性证明。

\subsection{Gram 矩阵、半正定规划与精确代数证书}
\label{b3:subsec:G05:02}

\begin{proposition}[平方和的 Gram 矩阵]\label{b3:prop:sos-gram-matrix}
设 \(n\ge1\)、\(d\ge0\) 为整数，\(f\in\mathbb R[x_1,\ldots,x_n]\) 的次数不超过 \(2d\)，\(v_d\) 为所有次数不超过 \(d\) 的单项式组成的列向量。则 \(f\) 为平方和，当且仅当存在实对称半正定矩阵 \(Q\) 使
\[
f=v_d^{\mathsf T}Qv_d.
\]
\end{proposition}
\begin{proof}
若 \(f=\sum_jq_j^2\)，最高次平方和不能抵消，故可取 \(\deg q_j\le d\)。写 \(q_j=c_j^{\mathsf T}v_d\)，则 \(Q=\sum_jc_jc_j^{\mathsf T}\) 半正定。反之，第二卷第10.2节的实二次型对角化给出 \(Q=C^{\mathsf T}C\)，其每行作用于 \(v_d\) 就得到所需平方。这里用的是半正定性；仅有实对称性会产生正负平方之差。
\end{proof}

展开 \(v_d^{\mathsf T}Qv_d=f\) 是关于 \(Q\) 各项的线性方程，加上 \(Q\succeq0\)，形成\term{半正定规划}[semidefinite programming]的可行性问题。这把第二卷的二次型用于更高次多项式；所增加的是单项式向量，而不是改变二次型定理。相关解释见\cite[Sections 7.1 and 7.3]{Sturmfels2002PolynomialEquations}。

\begin{example}\label{b3:exm:exact-gram}
令 \(v=(1,x,y)^{\mathsf T}\)，取
\[
Q=\begin{pmatrix}1&0&0\\0&1&-1\\0&-1&1\end{pmatrix}.
\]
则 \(v^{\mathsf T}Qv=1+(x-y)^2\)。矩阵和分解均为有理数，可逐项核验；其特征值为 \(1,0,2\)，零特征值并不妨碍证书成立。若把零特征值以浮点计算得到一个很小的负数，便不能直接断言矩阵半正定，须重新取得精确分解或其他严格证明。
\end{example}

对于求最小值问题，找到 \(f-\lambda\) 的约束平方和证书即可证明 \(\lambda\) 是下界；找到一个可行点达到它才完成最优性证明。若 \(M(g)\) 满足 Archimedean 条件、\(K\ne\varnothing\)，则每个 \(\lambda<\min_Kf\) 都有 Putinar 证书。由于每份证书只含有限次数，逐步提高允许次数得到的最优下界收敛到真正最小值。这一推论没有声称某个预先固定次数就能达到最优值。

\subsection{半代数拓扑、Nash 几何、o-minimal 几何与矩问题}
\label{b3:subsec:G05:03}

半代数集合也有欧氏拓扑。在 \(\mathbb R\) 上，开半代数集上的光滑半代数函数称为\term{Nash 函数}[Nash function]；例如 \(\sqrt{1+x^2}\) 是 Nash 函数，其图像由 \(y^2=1+x^2,y>0\) 刻画。定义与光滑代数几何的关系可读\cite[Appendix B, Definitions B.2.3--B.2.4]{Mangolte2020RealAlgebraicVarieties}。它兼有代数约束与微分性质，不能把任意光滑函数都归入此类。

\ProseParagraphBreak
半代数集合在直线上是有限个点与区间的并：组成其定义的有限多个非零多项式只有有限多个实根，在这些根之间每个多项式符号固定。\term{o-minimal 结构}[o-minimal structure]把这种一维有限性作为公理：在有序域各有限次幂上指定含代数集合、对布尔运算、直积和投影封闭的集合族，并要求一维集合都是有限个点与区间的并。Tarski--Seidenberg 定理保证半代数集合构成一个例子。更广结构、分层与拓扑不变量的理论不在此展开。

\ProseParagraphBreak
平方和的对偶计算引出\term{矩问题}[moment problem]。给定实数族 \((y_\alpha)_{\alpha\in\mathbb N^n}\)，先定义线性泛函 \(L(x^\alpha)=y_\alpha\)，再问是否存在支于 \(K\) 的正测度 \(\mu\)，使 \(L(p)=\int_Kp\,\mathrm{d}\mu\)。必要条件 \(L(q^2)\ge0\) 等价于所有有限矩矩阵 \((y_{\alpha+\beta})\) 半正定；若支集还满足 \(g_i\ge0\)，则须有 \(L(g_iq^2)\ge0\)。这些条件显然必要，但仅检查一张有限矩阵不能自动产生测度。

\ProseParagraphBreak
\begin{theorem}[Archimedean 二次模的矩表示]\label{b3:thm:archimedean-moment-representation}
设 \(g_1,\ldots,g_s\in\mathbb R[x_1,\ldots,x_n]\)，\(K=\{x:g_i(x)\ge0\text{ 对所有 }i\}\)，且存在 \(N>0\) 使 \(N-\sum_i x_i^2\in M(g)\)。若实线性泛函 \(L:\mathbb R[x]\to\mathbb R\) 满足 \(L(1)=1\) 与 \(L(M(g))\subseteq[0,\infty)\)，则存在支于 \(K\) 的 Borel 概率测度 \(\mu\)，使 \(L(p)=\int_Kp\,\mathrm{d}\mu\) 对每个多项式 \(p\) 成立。
\end{theorem}
\begin{proofsketch}
先由\cref{b3:thm:compact-positivity-certificates}的 Putinar 部分得出：若 \(p\ge0\) 于 \(K\)，则对每个 \(\varepsilon>0\)，\(p+\varepsilon\in M(g)\)，故 \(L(p)+\varepsilon\ge0\)，令 \(\varepsilon\to0\) 得 \(L(p)\ge0\)。\(K\) 不能为空，否则 \(-1\) 也严格正于空集，正性定理给出 \(-1\in M(g)\)，与 \(L(1)=1\) 矛盾。

若 \(p\) 在 \(K\) 上为零，分别用于 \(p\) 与 \(-p\) 得 \(L(p)=0\)，所以 \(L\) 下降为 \(K\) 上多项式函数的线性泛函。由 \(\|p\|_K\pm p\ge0\)，有 \(|L(p)|\le\|p\|_K\)。多项式含常数并分离 \(K\) 的点；Stone--Weierstrass 定理使它们在 \(C(K,\mathbb R)\) 中一致稠密。这个界给出 \(L\) 的唯一连续延拓，非负连续函数用多项式一致逼近后仍有非负像。

最后，紧 Hausdorff 空间上的 Riesz 正泛函表示定理把这个连续正泛函写成对正 Borel 测度的积分；\(L(1)=1\) 使其总质量为一。这里省略 Stone--Weierstrass 与 Riesz 两个分析定理的证明，保留它们的具体适用条件；从 Putinar 证书到矩表示的上述推导见\cite[Section 1, Theorem 2; Section 3, proof of Theorem 2]{Schweighofer2005Optimization}。
\end{proofsketch}

矩表示把多项式的正性条件转为正泛函，再转为测度；只检查一张有限矩矩阵则尚未提供对所有多项式所需的正性。



\begin{chaptersummary}
多项式等式的实解带有序结构；投影因此产生不等式，平方和因此产生超出普通根理想的信息。实零点定理把这种逐点信息写成有限代数证书，但更简洁的正性表示需要紧性或 Archimedean 等额外假设。Motzkin 多项式说明非负性与多项式平方和之间确有差别。

\ProseParagraphBreak
计算时，一元消去多项式可以继续交给 Sturm 序列，二次型可以通过单项式向量继续用于高次平方和。二者都要求保留精确数据：根隔离要能判定零与正负，正性证书要能核验恒等式与半正定性。浮点计算有助于寻找候选，最后的代数结论仍须由这些精确条件支持。
\end{chaptersummary}

\BookExercises

\begin{exercise}\label{b3:ex:G-projection-hyperbola}
在实闭域 \(R\) 上求 \(\{(x,y):xy=1,\ y>0\}\) 的 \(x\) 投影，并求 \((xy-1)\cap R[x]\)。
\end{exercise}
\begin{solution}
由 \(x=1/y\) 得投影为 \(x>0\)，反之取 \(y=1/x\)。若 \(h(x)\in(xy-1)\cap R[x]\)，向 \(R(x)\) 代入 \(y=1/x\) 得 \(h(x)=0\)，所以交为零。普通消去理想不保留正号条件。
\end{solution}

\begin{exercise}\label{b3:ex:G-parametric-quadratic}
消去量词 \(\exists y\,(y^2+ay+b=0)\)，并消去 \(\forall y\,(y^2+ay+b>0)\)。
\end{exercise}
\begin{solution}
配方为 \((y+a/2)^2+b-a^2/4\)。第一个条件等价于 \(a^2-4b\ge0\)，第二个等价于 \(4b-a^2>0\)：必要性在 \(y=-a/2\) 检验，充分性由平方非负。实闭性用于第一式的平方根存在。
\end{solution}

\begin{exercise}\label{b3:ex:G-real-radical-powers}
计算 \(I=(x^2+y^4)\subseteq\mathbb R[x,y]\) 的实根理想，并直接给出 \(x,y\) 属于它的证据。
\end{exercise}
\begin{solution}
实零点只有原点，故实根理想为 \((x,y)\)。直接看，\(x^2+(y^2)^2\in I\) 迫使任何包含 \(I\) 的实理想包含 \(x,y^2\)，进而包含 \(y\)。另一方面 \((x,y)\) 自身为实理想，因为商为 \(\mathbb R\)。因此无需调用实零点定理也能完成计算。
\end{solution}

\begin{exercise}\label{b3:ex:G-real-prime-versus-prime}
证明 \((x^2+1)\subseteq\mathbb R[x]\) 是素理想而不是实理想；证明零理想在 \(\mathbb R[x]\) 中是实理想。
\end{exercise}
\begin{solution}
前者商环为 \(\mathbb C\)，故为素理想；但 \(x^2+1^2\) 属于该理想而 \(1\) 不属于。若 \(\sum_jp_j^2=0\) 为多项式恒等式，在每个实数上求值得各 \(p_j\) 恒为零；一元非零多项式不能有无限多个根，故零理想实。
\end{solution}

\begin{exercise}\label{b3:ex:G-univariate-two-squares}
将 \((x^2+1)(x^2+4)\) 写成两个有理系数多项式的平方和。
\end{exercise}
\begin{solution}
取乘平方和恒等式中的 \((u,v,s,t)=(x,1,x,2)\)，得到 \((x^2-2)^2+(3x)^2=x^4+5x^2+4\)。
\end{solution}

\begin{exercise}\label{b3:ex:G-motzkin-zeros}
求 Motzkin 多项式 \(M=x^4y^2+x^2y^4+1-3x^2y^2\) 的全部实零点。
\end{exercise}
\begin{solution}
算术平均与几何平均取等号当且仅当三个非负数相等；这里第三个为一，故 \(x^4y^2=x^2y^4=1\)。于是 \(x,y\ne0\)，两式相除得 \(x^2=y^2\)，再得 \(x^6=1\)。实数中 \(x^2=y^2=1\)，四个符号组合都给零点。
\end{solution}

\begin{exercise}\label{b3:ex:G-empty-circle-certificate}
为实方程 \(x^2+y^2+1=0\) 写出不可行性证书；再给出条件 \(x^2+y^2\le1,\ x\ge2\) 的不可行性证书。
\end{exercise}
\begin{solution}
第一式有 \(-1=-(x^2+y^2+1)+x^2+y^2\)。第二式令 \(g_1=1-x^2-y^2,g_2=x-2\)，则
\[
-3=g_1+(x-2)^2+4g_2+y^2.
\]
除以三，所有非恒等项的系数均为正实数平方或平方和，便得到 \(-1\in T(g_1,g_2)\)。
\end{solution}

\begin{exercise}\label{b3:ex:G-quadratic-module-ball}
在由 \(g_i=1-x_i^2\) 生成的二次模中，证明 Archimedean 球约束成立。
\end{exercise}
\begin{solution}
直接相加得 \(n-\sum_i x_i^2=\sum_i g_i\in M(g)\)。每个系数一都是平方，故这是精确证书；相应的可行集是闭立方体。
\end{solution}

\begin{exercise}\label{b3:ex:G-strict-positivity-warning}
解释为何对任意 \(\varepsilon>0\)，\(1-x^2+\varepsilon\) 属于 \(M((1-x^2)^3)\)，而 \(1-x^2\) 不属于。
\end{exercise}
\begin{solution}
正文已证明该二次模满足 Archimedean 条件，其可行集为 \([-1,1]\)。加 \(\varepsilon\) 后严格正，由 Putinar 定理属于二次模；不加时，用正文在 \(x=1\) 比较一阶导数的论证得到矛盾。这说明一列存在的正性证书不必在极限处保持同样的表示。
\end{solution}

\begin{exercise}\label{b3:ex:G-sturm-cubic}
用 Sturm 序列证明 \(x^3-x-1\) 只有一个实根，并把它隔离在 \((1,3/2)\)。
\end{exercise}
\begin{solution}
可用正比例简化后的序列 \(x^3-x-1,3x^2-1,2x+3,-1\)。在 \(-\infty\) 的符号为 \(-,+,-,-\)，变化二次；在 \(+\infty\) 为 \(+,+,+,-\)，变化一次，总实根数一。在 \(1\) 的符号为 \(-,+,+,-\)，变化二次；在 \(3/2\) 为 \(+,+,+,-\)，变化一次。因此该区间恰含唯一实根。端点均非根。
\end{solution}

\begin{exercise}\label{b3:ex:G-repeated-roots}
对 \(p=(x-1)^4(x+2)^2(x^2+1)\) 说明 Sturm 计数前如何处理重根，并列出所有实根及重数。
\end{exercise}
\begin{solution}
\(\gcd(p,p')=(x-1)^3(x+2)\)，除后得到平方自由部分 \((x-1)(x+2)(x^2+1)\)。它的实根为 \(1,-2\)，各按不同根计数一次；原分解记录的重数分别为四、二。根计数与重数恢复是两个相关而不同的输出。
\end{solution}

\begin{exercise}\label{b3:ex:G-eliminate-and-isolate}
精确求解实方程组 \(y=x^2,\ y^2-5y+4=0\)，说明消去与 Sturm 隔离各完成什么工作。
\end{exercise}
\begin{solution}
代入第一式得到 \(x^4-5x^2+4=0\)，正文的四个隔离区间各含一根。进一步分解得 \(x=\pm1,\pm2\)，回代得解 \((\pm1,1),(\pm2,4)\)。消去得到候选一元方程，隔离保证实根的个数与位置；最后回代验证原方程。一般消去所得候选未必全可提升，不能省略最后一步。
\end{solution}

\begin{exercise}\label{b3:ex:G-gram-exact}
为 \(f=x^4+2x^2y^2+y^4+1\) 在 \(v=(1,x^2,xy,y^2)^{\mathsf T}\) 下写出一个精确 Gram 矩阵。
\end{exercise}
\begin{solution}
由于 \(f=1+(x^2+y^2)^2\)，可取
\[
Q=\begin{pmatrix}1&0&0&0\\0&1&0&1\\0&0&0&0\\0&1&0&1\end{pmatrix}.
\]
它是向量 \((1,0,0,0)^{\mathsf T}\) 与 \((0,1,0,1)^{\mathsf T}\) 的外积之和，因而半正定。逐项相乘即可验证恒等式。
\end{solution}

\begin{exercise}\label{b3:ex:G-certified-minimum}
在 \(x^2+y^2\le1\) 上求 \(f=(x-y)^2+x^2+y^2\) 的最小值，并同时给出下界证书和取等点。
\end{exercise}
\begin{solution}
\(f=(x-y)^2+x^2+y^2\) 本身为平方和，故下界零在整个平面上成立。点 \((0,0)\) 满足约束并取零，故最小值为零。只有下界恒等式而没有取等点时，还不能宣布该下界最优。
\end{solution}

\begin{exercise}\label{b3:ex:G-moment-two-points}
令 \(\mu=\frac12\delta_{-1}+\frac12\delta_1\)。求前五个矩，并检验次数不超过二的矩矩阵半正定。
\end{exercise}
\begin{solution}
\(y_0=1,y_1=0,y_2=1,y_3=0,y_4=1\)，故矩矩阵为
\[
\begin{pmatrix}1&0&1\\0&1&0\\1&0&1\end{pmatrix}.
\]
对系数向量 \((a,b,c)\)，其二次型为 \((a+c)^2+b^2\ge0\)，也等于 \(\int(a+bx+cx^2)^2\,\mathrm{d}\mu\)。这里测度已经给出，所以必要条件的验证不涉及一般矩表示存在性。
\end{solution}
